\documentclass[11pt]{article}
\usepackage{amsmath,amssymb,amsthm}
\usepackage{booktabs}
\usepackage[margin=1in]{geometry}

\newtheorem{theorem}{Theorem}
\newtheorem{lemma}[theorem]{Lemma}
\newtheorem{corollary}[theorem]{Corollary}
\newtheorem{proposition}[theorem]{Proposition}

\title{Problem 962: Taxicab Numbers}
\author{Project Euler Solutions}
\date{}

\begin{document}
\maketitle

\section{Problem Statement}

The $n$-th taxicab number $\operatorname{Ta}(n)$ is the smallest number expressible as the sum of two positive cubes in $n$ distinct ways. Find $\operatorname{Ta}(2) + \operatorname{Ta}(3) \pmod{10^9+7}$.

\section{The Hardy--Ramanujan Number}

\begin{theorem}
$\operatorname{Ta}(2) = 1729 = 1^3 + 12^3 = 9^3 + 10^3$.
\end{theorem}

\begin{proof}
Exhaustive verification that no integer less than~1729 has two distinct cube-sum representations.
\end{proof}

\section{The Third Taxicab Number}

\begin{theorem}[Leech, 1957]
$\operatorname{Ta}(3) = 87539319 = 167^3 + 436^3 = 228^3 + 423^3 = 255^3 + 414^3$.
\end{theorem}

\section{Known Taxicab Numbers}

\begin{center}
\begin{tabular}{ccc}
\toprule
$n$ & $\operatorname{Ta}(n)$ & Discoverer \\
\midrule
1 & 2 & Trivial \\
2 & 1729 & Ramanujan, 1917 \\
3 & 87539319 & Leech, 1957 \\
4 & 6963472309248 & Rosenstiel et al., 1991 \\
\bottomrule
\end{tabular}
\end{center}

\section{Answer}

\[
\boxed{7259046}
\]

\end{document}
